\chapter{Server Communication}
\label{ch:server-comm}

All communication between the Kiss frontend and backend uses JSON-RPC via the \texttt{Server}\index{Server.js} class.  This chapter covers making calls, handling responses, managing sessions, uploading files, and downloading reports.

\section{Server.call()}
\label{sec:sc-call}

\texttt{Server.call()} is the primary communication method.  It sends a JSON-RPC request to the backend and returns a Promise with the response:

\begin{lstlisting}[style=javascript]
const res = await Server.call("services/Crud", "getRecords", {
    departmentId: 42,
    activeOnly: true
});
if (res._Success) {
    grid.addRecords(res.rows);
}
\end{lstlisting}

\noindent
The three required parameters are:

\begin{enumerate}
  \item \textbf{Service path} --- the path to the backend service class (e.g., \texttt{"services/Crud"}).
  \item \textbf{Method name} --- the method to invoke on the service class.
  \item \textbf{Input data} --- a JavaScript object sent as the \texttt{injson} parameter to the backend method.  This parameter is optional; omit it if the method requires no input.
\end{enumerate}

\noindent
\texttt{Server.call()} automatically:
\begin{itemize}
  \item Includes the session UUID in every request.
  \item Shows a wait cursor while the request is in progress.
  \item Handles authentication failures (error code 2) by redirecting to login.
  \item Retries once on transient network errors.
  \item Displays error messages for server-side exceptions.
\end{itemize}

\section{Request/Response Format}
\label{sec:sc-format}

Behind the scenes, \texttt{Server.call()} sends an HTTP POST to the \texttt{/rest} endpoint with a JSON body:

\begin{lstlisting}[style=javascript]
// Outgoing request
{
    "_class": "services/Crud",
    "_method": "getRecords",
    "_uuid": "a1b2c3d4-e5f6-...",
    "departmentId": 42,
    "activeOnly": true
}

// Successful response
{
    "_Success": true,
    "_ErrorCode": 0,
    "rows": [ ... ]
}

// Error response
{
    "_Success": false,
    "_ErrorCode": -1,
    "_ErrorMessage": "Department not found"
}
\end{lstlisting}

\noindent
The \texttt{\_Success} field is always present.  On error, \texttt{\_ErrorMessage} contains the error text and \texttt{\_ErrorCode} indicates the type:

\begin{longtable}{cl}
\toprule
\textbf{Code} & \textbf{Meaning} \\
\midrule
\endhead
0 & Success \\
2 & Authentication failure (session expired or invalid) \\
$-1$ & Server error (exception in service method) \\
\bottomrule
\end{longtable}

\section{Error Handling}
\label{sec:sc-errors}

When the backend throws an exception, \texttt{Server.call()} displays the error message to the user automatically.  The callback still executes, but \texttt{res.\_Success} is \texttt{false}:

\begin{lstlisting}[style=javascript]
const res = await Server.call(WS, 'save', data);
if (res._Success) {
    // handle success
    Utils.showMessage('Info', 'Saved successfully');
} else {
    // The error message has already been shown to the user.
    // This branch is optional --- use it only if you need
    // custom error recovery.
}
\end{lstlisting}

\noindent
Different backend exception types produce different behaviors:

\begin{description}[style=nextline]
  \item[\texttt{UserException}] The message is displayed to the user.  No stack trace is logged on the server.
  \item[\texttt{ServerException}] A generic error message is shown.  The full stack trace is logged on the server.
  \item[\texttt{LogException}] The message is displayed to the user \emph{and} the full stack trace is logged on the server.
\end{description}

\section{Quiet Calls}
\label{sec:sc-quiet}

\texttt{Server.callQuiet()} is a variant of \texttt{Server.call()} for lightweight background requests --- type-ahead searches, autocomplete lookups, polling, and similar work where the screen should not appear blocked and where a transport failure should be handled silently rather than interrupting the user.  It takes the same arguments as \texttt{Server.call()} and returns the same Promise-wrapped response:

\begin{lstlisting}[style=javascript]
const res = await Server.callQuiet("services/Customers",
                                   "search", {q: text});
if (res._Success)
    showResults(res.customers);
\end{lstlisting}

\noindent
It differs from \texttt{Server.call()} in two ways:

\begin{itemize}
  \item It does \emph{not} change the global wait cursor or busy count, so a rapid sequence of keystroke-driven calls does not make the whole application flicker as busy.  The screen owns whatever loading indication it wants to show.
  \item A transport or JSON-parse failure does \emph{not} raise the framework error popup.  Instead the caller receives a \texttt{\{\_Success: false\}} result to handle quietly.
\end{itemize}

\noindent
In every other respect it behaves like \texttt{Server.call()}: the session UUID is attached automatically, and an authentication failure (error code 2) still shows the standard message and routes to the login screen.

\begin{tipbox}
Use \texttt{callQuiet()} for background work whose result the user did not explicitly request, and \texttt{call()} for actions the user is waiting on (saving a form, running a report).  The \texttt{<search-input>} component (Chapter~\ref{ch:components}) is designed to pair with \texttt{callQuiet()}.
\end{tipbox}

\section{Parallel Calls}
\label{sec:sc-parallel}

When a screen needs data from multiple independent services, use \texttt{Server.callAll()} to execute them in parallel:

\begin{lstlisting}[style=javascript]
const error = await Server.callAll(
    [
        Server.call(WS, 'getDepartments'),
        Server.call(WS, 'getStatuses')
    ],
    function (deptRes) {
        // process department data
        for (let dept of deptRes.departments)
            $$('department').add(dept.id, dept.name);
    },
    function (statusRes) {
        // process status data
        for (let s of statusRes.statuses)
            $$('status').add(s.id, s.name);
    }
);
if (!error) {
    // all calls succeeded
}
\end{lstlisting}

\noindent
\texttt{callAll()} waits for all promises to complete before calling any handlers.  It returns \texttt{true} if any call failed, \texttt{false} if all succeeded.

\section{Session Management}
\label{sec:sc-sessions}

After login, the backend assigns a session token (a UUID).  \texttt{Server.js} stores this token and includes it automatically in every subsequent request as the \texttt{\_uuid} field.

The session expires after a configurable period of inactivity (set by \texttt{UserInactiveSeconds} in \texttt{application.ini}).  When a request arrives with an expired UUID, the backend returns error code 2 and \texttt{Server.call()} redirects to the login screen.

\subsection{Inactivity Timeout (Frontend)}

The frontend can also enforce an inactivity timeout independently:

\begin{lstlisting}[style=javascript]
Server.setMaxInactivitySeconds(900);   // 15 minutes
Server.setMaxInactivityMinutes(15);    // equivalent
Server.setMaxInactivityHours(1);       // 1 hour
\end{lstlisting}

\noindent
The frontend timeout runs a \texttt{checkTime()} method that compares the elapsed time since the last user interaction.  If the timeout is exceeded, the user is logged out.

\section{Logging Out}
\label{sec:sc-logout}

\begin{lstlisting}[style=javascript]
Server.logout();
\end{lstlisting}

\noindent
\texttt{Server.logout()} performs a complete cleanup:

\begin{enumerate}
  \item Calls the backend \texttt{Logout} service to invalidate the session.
  \item Clears the stored UUID (and all other persisted client state).
  \item Disables back-button protection.
  \item Calls \texttt{Utils.cleanup()} to release frontend resources.
  \item Performs a full page reload of a fresh \texttt{index.html}, landing on the login route (\texttt{Router.loginHash()}).
\end{enumerate}

\noindent
The reload in the final step re-boots the tab on the latest cache-busted code, so a tab that stayed open across a redeploy picks up the new release the next time its session ends (Chapter~\ref{ch:caching-security}).

\section{File Uploads}
\label{sec:sc-uploads}

File uploads use \texttt{Server.fileUploadSend()}, which sends files as multipart/form-data:

\begin{lstlisting}[style=javascript]
$$('upload-btn').onclick(async () => {
    const count = $$('file-ctl').numberOfUploadFiles();
    if (count === 0) {
        Utils.showMessage('Error', 'Please select a file');
        return;
    }
    const fd = $$('file-ctl').getFormData();
    const res = await Server.fileUploadSend(
        WS, 'upload', fd, {},
        'Uploading...', 'Upload complete');
});
\end{lstlisting}

\noindent
The parameters are:

\begin{enumerate}
  \item \textbf{Service path} and \textbf{method name} --- same as \texttt{Server.call()}.
  \item \textbf{FormData} --- from the \texttt{<file-upload>} component's \texttt{getFormData()}.
  \item \textbf{Additional data} --- a JavaScript object with extra fields to send alongside the files.
  \item \textbf{Wait message} --- shown during upload.
  \item \textbf{Success message} --- shown after successful upload.
\end{enumerate}

\noindent
On the backend, the service method accesses uploaded files through the \texttt{servlet} parameter:

\begin{lstlisting}[style=groovy]
int count = servlet.getUploadFileCount()
String name = servlet.getUploadFileName(0)
BufferedInputStream is =
    servlet.getUploadBufferedInputStream(0)
\end{lstlisting}

\section{File Downloads}
\label{sec:sc-downloads}

For reports and exports, the backend generates a temporary file and returns its URL.  The frontend opens the URL to trigger the browser's download or display behavior:

\begin{lstlisting}[style=javascript]
// Backend generates a PDF report and returns the URL
const res = await Server.call(WS, 'runReport');
if (res._Success)
    Utils.showReport(res.reportUrl);
\end{lstlisting}

\begin{lstlisting}[style=groovy]
// Backend service method
Groff rpt = new Groff("Report-", "My Report", false)
// ... build report ...
outjson.put("reportUrl", rpt.process())
\end{lstlisting}

\noindent
Temporary report files are automatically cleaned up by the framework after approximately one day.

\section{Binary Data}
\label{sec:sc-binary}

For retrieving binary data (e.g., images), use \texttt{Server.binaryCall()}:

\begin{lstlisting}[style=javascript]
const res = await Server.binaryCall(WS, 'getImage',
    { id: imageId });
if (res._Success) {
    // res._data contains the binary content as an ArrayBuffer
    const blob = new Blob([res._data]);
    const url = URL.createObjectURL(blob);
    // use the URL in an img tag or picture component
}
\end{lstlisting}

\noindent
\texttt{binaryCall()} returns the response as an ArrayBuffer.  JSON metadata and binary content are separated by a delimiter character, and the parsing is handled automatically.

\section{A Complete Screen Example}
\label{sec:sc-complete-example}

The following example ties together everything in this chapter: loading data, populating controls, saving with validation, and refreshing:

\begin{lstlisting}[style=javascript]
'use strict';

(async function () {
    const WS = 'services/Crud';

    // --- Grid setup ---
    const columnDefs = [
        {headerName: 'Last Name', field: 'lastName',
            width: 33},
        {headerName: 'First Name', field: 'firstName',
            width: 33},
        {headerName: 'Phone', field: 'phoneNumber', width: 33}
    ];
    const grid = new AGGrid('grid', columnDefs, 'id');
    grid.show();

    // --- Load data ---
    async function updateGrid() {
        grid.clear();
        $$('edit').disable();
        $$('delete').disable();
        const res = await Server.call(WS, 'getRecords');
        if (res._Success)
            grid.addRecords(res.rows);
    }
    updateGrid();

    // --- Grid selection ---
    grid.setOnSelectionChanged((rows) => {
        $$('edit').enable(rows);
        $$('delete').enable(rows);
    });

    // --- Add new record ---
    $$('new').onclick(() => {
        $$('erp-first-name').clear();
        $$('erp-last-name').clear();
        $$('erp-phone-number').clear();
        Utils.popup_open('edit-record-popup',
            'erp-first-name');

        $$('erp-ok').onclick(async () => {
            if ($$('erp-first-name').isError('First Name'))
                return;
            if ($$('erp-last-name').isError('Last Name'))
                return;
            const data = {
                firstName: $$('erp-first-name').getValue(),
                lastName: $$('erp-last-name').getValue(),
                phoneNumber: $$('erp-phone-number').getValue()
            };
            const res = await Server.call(WS, 'addRecord',
                data);
            if (res._Success) {
                Utils.popup_close();
                updateGrid();
            }
        });

        $$('erp-cancel').onclick(() => Utils.popup_close());
    });

    // --- Delete with confirmation ---
    $$('delete').onclick(() => {
        Utils.yesNo('Confirm',
            'Delete the selected record?', async () => {
            const row = grid.getSelectedRow();
            const res = await Server.call(WS, 'deleteRecord',
                { id: row.id });
            if (res._Success) updateGrid();
        });
    });

    // --- PDF report ---
    $$('report').onclick(async () => {
        const res = await Server.call(WS, 'runReport');
        if (res._Success)
            Utils.showReport(res.reportUrl);
    });

    // --- CSV export ---
    $$('export').onclick(async () => {
        const res = await Server.call(WS, 'runExport');
        if (res._Success)
            Utils.showReport(res.exportUrl);
    });
})();
\end{lstlisting}

\noindent
This is adapted from the actual CRUD example that ships with Kiss.  It demonstrates the standard patterns: load data into a grid, enable/disable buttons based on selection, open popups for editing, validate before saving, and confirm before deleting.
